\documentclass[11pt]{ctexart}
\usepackage[a4paper,left=4.8cm,right=4.8cm,top=4cm,bottom=4cm,headheight=14pt]{geometry}
\usepackage{amsmath}
\usepackage{amssymb}
\usepackage{tikzflora}

\title{花间数学}
\author{XeLaTeX 效果示例}
\date{}

\florasetup{
  layout=left,
  seed=7,
  colourway=spring,
  visitors=0.5,
  density=0.9,
  giltdpi=150
}

\begin{document}

\thispagestyle{empty}
\begin{center}
  {\Huge\bfseries 花间数学\par}
  \vspace{1em}
  {\large 一份使用 XeLaTeX 编译的水彩页边示例\par}
\end{center}
\vspace{2em}

\section{函数与图像}
设 $f(x)=x^2-2x-3$。将其配方，得到
\begin{equation}
  f(x)=(x-1)^2-4.
\end{equation}
因此抛物线的顶点为 $(1,-4)$，对称轴为 $x=1$。令 $f(x)=0$，可得两个零点：
\[
  x_1=-1,\qquad x_2=3.
\]

\section{单调性}
\begin{itemize}
  \item 当 $x<1$ 时，函数单调递减；
  \item 当 $x>1$ 时，函数单调递增；
  \item 函数的最小值为 $-4$。
\end{itemize}

\newpage
\florathispage{layout=right, colourway=autumn, visitors=0.8}
\section{三角函数}
对于任意实数 $x$，都有
\[
  \sin^2 x+\cos^2 x=1.
\]
由此可以进一步得到
\[
  1+\tan^2 x=\sec^2 x,
  \qquad x\neq \frac{\pi}{2}+k\pi,\ k\in\mathbb{Z}.
\]

\section{基本不等式}
若 $a>0$，$b>0$，则
\begin{equation}
  \frac{a+b}{2}\geq \sqrt{ab},
\end{equation}
等号成立当且仅当 $a=b$。这个结论常用于求最值。

\newpage
\florathispage{layout=centered, colourway=winter, visitors=0.2}
\section{数列}
等差数列 $\{a_n\}$ 的通项公式为
\[
  a_n=a_1+(n-1)d,
\]
前 $n$ 项和为
\[
  S_n=\frac{n(a_1+a_n)}{2}
     =na_1+\frac{n(n-1)}2d.
\]

等比数列 $\{b_n\}$ 的通项公式为 $b_n=b_1q^{n-1}$。当 $q\neq1$ 时，
\[
  T_n=\frac{b_1(1-q^n)}{1-q}.
\]

\newpage
\thispagestyle{empty}
\florathispage{layout=surround, colourway=summer, density=0.75}
\begin{center}
  {\Huge\bfseries 夏日小结\par}
\end{center}
\vspace{2em}

这一页使用环绕布局。植物由外部 \texttt{texlua} 进程生成单页 PDF，再由 XeLaTeX 放入页面背景，正文、中文字体和数学公式仍由 XeLaTeX 排版。

\begin{center}
\Large
$\displaystyle e^{i\pi}+1=0$
\end{center}

\vfill
\begin{center}
\small 固定随机种子后，同一页面可以重复生成相同的植物。
\end{center}

\end{document}
